% !PRINTED 389
\weberpage{389}
%<<CONT
fraction development finally obtain the same closing numbers.
If, then, both continued fractions are periodic, they must belong
to the same period.

If for a given discriminant $D$ one has developed the complete
system of reduced numbers according to the rules given at the
beginning of the section, it is easy to arrange these numbers
in periods and to find the partial denominators $\alpha$ of the continued
fraction. Namely, let
\weq{12}{\omega = \frac{\sqrt{D}+b}{2c} = \frac{2a}{\sqrt{D}-b}}
be one of these numbers and
\weq{13}{\omega_1 = \frac{\sqrt{D}+b_1}{2c_1} = \frac{2a_1}{\sqrt{D}-b_1}}
the number immediately following it in the period, so that
\weq{14}{\omega = \alpha + \frac{1}{\omega_1}, \quad \omega_1 = \frac{1}{\omega-\alpha}}
where $\alpha$ denotes the greatest integer contained in $\omega$.
By (12) and (14) we then have
\[
\omega_1 = \frac{2c}{\sqrt{D}+b-2c\alpha},
\]
and comparison with (13) gives
\weq{15}{a_1 = c, \quad b_1 = 2c\alpha - b.}
If, conversely, two of the reduced numbers $\{a, b, c\}$, $\{a_1, b_1, c_1\}$
belonging to $D$ stand in the relation represented by (15),
in which $\alpha$ is any integer, then the second immediately
follows the first in the period. For from (15)
follows (14), and hence by 2. $\alpha$ must be the greatest number
contained in $\omega$.

The number $\omega_1$ is therefore completely and uniquely determined
from $\omega$ by the conditions (15).

7. \gesperrt{One therefore arranges the numbers $\{a, b, c\}$ from left
to right in such a way that the last number $c$
of the preceding one becomes at the same time the first number $a_1$ of the
following, and that the sum of the two
middle numbers $b + b_1$ is divisible by $2c$.
This arrangement is possible in only one way,
and the quotient $b + b_1 : 2c$ is the number $\alpha$ that
appears as partial denominator in the continued fraction.}
